% !TEX program = xelatex
\documentclass[11pt,a4paper]{article}

% === CJK Support ===
\usepackage[UTF8]{ctex}
\usepackage{fontspec}
\setCJKmainfont{SimSun}[BoldFont=SimHei,ItalicFont=KaiTi]
\setmainfont{Times New Roman}

% === Math ===
\usepackage{amsmath,amssymb,amsthm}
\usepackage{mathtools}
\usepackage{bm}

% === Theorem Environments ===
\newtheorem{theorem}{定理}[section]
\newtheorem{lemma}[theorem]{引理}
\newtheorem{corollary}[theorem]{推论}
\newtheorem{proposition}[theorem]{Proposition}
\newtheorem{definition}[theorem]{定义}
\newtheorem{remark}[theorem]{注记}

% === Honest Strike标注 ===
\newtheorem{honestattack}{Honest Strike}
\renewcommand{\thehonestattack}{\textcolor{red}{\textbf{[!] Strike \arabic{honestattack}}}}
\newenvironment{attackbox}{%
  \begin{center}
  \color{red}
  \begin{minipage}{0.92\textwidth}
  \begin{honestattack}
}{%
  \end{honestattack}
  \end{minipage}
  \end{center}
}

% === Misc ===
\usepackage{hyperref}
\usepackage{geometry}
\geometry{margin=2.5cm}
\usepackage{graphicx}
\usepackage{enumitem}
\usepackage{booktabs}
\usepackage{xcolor}
\usepackage{tikz}

\hypersetup{
  colorlinks=true,
  linkcolor=blue,
  citecolor=blue,
  urlcolor=blue
}

% === Macros ===
\newcommand{\R}{\mathbb{R}}
\newcommand{\E}{\mathbb{E}}
\newcommand{\Var}{\mathrm{Var}}
\newcommand{\Cov}{\mathrm{Cov}}
\newcommand{\indep}{\perp\!\!\!\perp}
\newcommand{\logit}{\mathrm{logit}}
\newcommand{\sgn}{\mathrm{sgn}}
\newcommand{\err}{\mathrm{err}}
\renewcommand{\P}{\mathbb{P}}
\newcommand{\one}{\mathbf{1}}
\newcommand{\wv}{\mathbf{w}}
\newcommand{\xv}{\mathbf{x}}
\newcommand{\pv}{\mathbf{p}}
\newcommand{\Sigmav}{\mathbf{\Sigma}}
\newcommand{\Corr}{\mathrm{Corr}}

\title{\textbf{Collective Intelligence:  Condorcet 陪审团定理的现代形式化}\\[5pt]
       \large 从 SCX Theorem~1 出发的四重推广}
\author{SCX \\[3pt] \small 独立研究}
\date{2026年6月}

\begin{document}

\maketitle

\begin{abstract}
Condorcet 陪审团定理（CJT）奠定了集体智能的数学基础：
当每位独立专家正确概率 $p > 1/2$ 时，简单多数投票的准确率随群体规模增大而趋于~1。
本文从作者此前建立的 SCX Theorem~1（同质独立专家的最优多数界）出发，
给出四个方向的形式化推广：
(1)~\textbf{异质专家最优加权定理}——专家准确率不同时，对数几率加权达到贝叶斯最优；
(2)~\textbf{多样性-规模权衡定理}——$M$ 个相关专家与 $M'$ 个独立专家的效率awaiting价条件；
(3)~\textbf{相关性结构下的共识下界}——已知相关矩阵时多数投票错误的非平凡下界；
(4)~\textbf{策略性投票的审计检测界}——审计者以有限样本检测策略性谎报的统计极限。
每个定理均附形式Proof与数值Example，
并以\textcolor{red}{\textbf{Honest Strike}}标注其边界条件、
隐藏假设与不能成立的情形。

\medskip
\noindent\textbf{Keywords：} 集体智能；Condorcet 陪审团定理；异质加权投票；多样性-规模权衡；
相关性结构；策略性投票审计；统计学习理论
\end{abstract}

\tableofcontents

% ===================================================================
\section{引言}
% ===================================================================

Condorcet 陪审团定理是现代集体决策理论的基石。
考虑 $n$ 位专家对二元假设 $\theta \in \{-1,+1\}$ 各自投票~$X_i \in \{-1,+1\}$。
经典定理断言：若每位专家的判断 $X_i$ 条件独立于 $\theta$，且正确概率
$\P(X_i = \theta \mid \theta) = p > 1/2$，则简单多数投票
$\hat{\theta} = \sgn\bigl(\sum_{i=1}^n X_i\bigr)$ 的错误概率随~$n$ 指数Decays。

然而，经典表述存在三个关键的不现实假设：
\begin{enumerate}[label=(\roman*)]
  \item \textbf{同质性}——所有专家共享相同准确率~$p$；
  \item \textbf{独立性}——专家判断彼此条件独立；
  \item \textbf{诚实性}——专家总是如实报告其判断。
\end{enumerate}

本文从一条基线定理出发，逐层放松上述三条假设。

\begin{theorem}[SCX Theorem~1 — 同质独立多数投票界]
\label{thm:scx1}
设 $X_1,\ldots,X_n \in \{-1,+1\}$ 关于 $\theta \in \{-1,+1\}$ 条件独立，
且 $\P(X_i = \theta \mid \theta) = p > 1/2$ 对所有 $i$ 成立。
令 $\hat{\theta}_n = \sgn(S_n)$，其中 $S_n = \sum_{i=1}^n X_i$。
则对任意 $n \geq 1$，
\begin{equation}
  \P(\hat{\theta}_n \neq \theta) \leq \exp\!\bigl(-2n(p-\tfrac{1}{2})^2\bigr).
  \label{eq:hoeffding}
\end{equation}
更精确地，当 $n$ 为奇数时，
\begin{equation}
  \P(\hat{\theta}_n \neq \theta) = \sum_{k=0}^{\lfloor n/2 \rfloor} \binom{n}{k} p^{k}(1-p)^{n-k}
  \leq \frac{1}{2}\bigl[2\sqrt{p(1-p)}\bigr]^n.
  \label{eq:exact}
\end{equation}
\end{theorem}

\begin{proof}
式 (\ref{eq:hoeffding}) 是 $S_n/n$ 作为 $p$ 的样本平均的 Hoeffding 界直接结果；
式 (\ref{eq:exact}) 来自二项分布的精确尾部与 Chernoff--Bhattacharyya 界。
\end{proof}

\begin{attackbox}
  \textbf{SCX Theorem~1 的脆弱性}。
  界~(\ref{eq:hoeffding}) 的指数速率 $2(p-1/2)^2$ 是紧的（达到 Cram\'er--Chernoff 大偏差速率），
  但其前提——条件独立与同质准确率——在真实集体决策中几乎从未成立。
  即使微小违背独立性（如专家间的信息重叠），也将系统性地破坏指数Decays。
  本文的后续定理即围绕这一裂缝展开。
\end{attackbox}

% ===================================================================
\section{预备知识}
% ===================================================================

\subsection{记号与设定}

考虑二元Status空间 $\Theta = \{-1,+1\}$，先验 $\P(\theta = +1) = \pi \in (0,1)$。
$n$ 位专家各自发出信号 $X_i \in \{-1,+1\}$。
记向量 $\xv = (X_1,\ldots,X_n)^\top$。

\begin{definition}[信号剖面]
专家 $i$ 的\emph{信号剖面}定义为她的第一类错误率 $\alpha_i = \P(X_i = -1 \mid \theta = +1)$
与第二类错误率 $\beta_i = \P(X_i = +1 \mid \theta = -1)$。
准确率 $p_i = \P(X_i = \theta \mid \theta)$ 满足 $p_i = 1 - \frac{\alpha_i + \beta_i}{2}$（当先验均awaiting时）。
\end{definition}

\begin{definition}[对数几率权重]
对于错误率 $(\alpha_i,\beta_i)$ 的专家，其\emph{对数几率权重}定义为
\begin{equation}
  w_i^* = \log\frac{1-\alpha_i}{\beta_i} - \log\frac{\alpha_i}{1-\beta_i}
        = \log\frac{(1-\alpha_i)(1-\beta_i)}{\alpha_i\beta_i}.
  \label{eq:logoddsweight}
\end{equation}
当 $\alpha_i = \beta_i = 1-p_i$（对称错误），则 $w_i^* = 2\log\frac{p_i}{1-p_i} = 2\,\logit(p_i)$。
\end{definition}

\subsection{决策规则族}

本文考虑的加权多数决策规则为：
\begin{equation}
  \hat{\theta}_{\wv}(\xv) = \sgn\!\left(\sum_{i=1}^n w_i X_i\right),
  \label{eq:weightedrule}
\end{equation}
其中 $\wv = (w_1,\ldots,w_n)^\top \in \R_{\geq 0}^n$ 为权重向量。
$w_i = 1$ 对所有 $i$ 退化为简单多数投票。

% ===================================================================
\section{异质专家最优加权定理}
% ===================================================================

本节回答：\textbf{当专家准确率不同时，如何分配投票权重以最大化集体正确概率？}

\subsection{二元决策的贝叶斯最优权重}

\begin{theorem}[异质专家最优加权定理]
\label{thm:optimalweight}
设 $X_1,\ldots,X_n$ 关于 $\theta$ 条件独立，专家 $i$ 的（可能非对称）错误率为 $(\alpha_i,\beta_i)$。
则规则族 (\ref{eq:weightedrule}) 中使错误概率最小化的权重为
\begin{equation}
  w_i^* = \log\frac{1-\alpha_i}{\beta_i} + \log\frac{1-\beta_i}{\alpha_i}
        = \log\frac{(1-\alpha_i)(1-\beta_i)}{\alpha_i\beta_i},
  \label{eq:optw}
\end{equation}
且该权重awaiting于对数后验几率比：
\begin{equation}
  \sum_{i=1}^n w_i^* X_i = \log\frac{\P(\theta=+1 \mid \xv)}{\P(\theta=-1 \mid \xv)}
                         - \log\frac{\pi}{1-\pi}.
  \label{eq:posterior}
\end{equation}
\end{theorem}

\begin{proof}
在条件独立假设下，似然比分解为
\[
  \frac{\P(\xv \mid \theta=+1)}{\P(\xv \mid \theta=-1)}
  = \prod_{i: X_i=+1} \frac{1-\alpha_i}{\beta_i}
    \prod_{i: X_i=-1} \frac{\alpha_i}{1-\beta_i}.
\]
取对数：
\[
  \log\frac{\P(\xv \mid \theta=+1)}{\P(\xv \mid \theta=-1)}
  = \sum_{i=1}^n \left[
      \one\{X_i=+1\} \log\frac{1-\alpha_i}{\beta_i}
    + \one\{X_i=-1\} \log\frac{\alpha_i}{1-\beta_i}
    \right].
\]
注意 $\one\{X_i=+1\} = \frac{1+X_i}{2}$，$\one\{X_i=-1\} = \frac{1-X_i}{2}$。
代入并化简即得加权和形式 $\sum_i w_i^* X_i + \text{const}$，
其中 $w_i^*$ 如 (\ref{eq:optw}) 所示。
贝叶斯最优决策规则为 $\sgn\bigl(\log\frac{\P(\theta=+1|\xv)}{\P(\theta=-1|\xv)}\bigr)$，
它awaiting价于以 $w_i^*$ 加权的多数规则（至多差一个与数据无关的平移）。
\end{proof}

\begin{corollary}[对称情形的经典Conclusion]
若专家 $i$ 的错误对称（$\alpha_i = \beta_i = \varepsilon_i$），则
$w_i^* = 2\log\frac{1-\varepsilon_i}{\varepsilon_i}
       = 2\log\frac{p_i}{1-p_i}$。
即\textbf{每位专家的最优权重与其对数几率成正比}。
\end{corollary}

\begin{corollary}[同质退化为 SCX Theorem~1]
若所有 $\varepsilon_i = \varepsilon$，则所有权重相awaiting，规则退化为简单多数投票，
错误界退化为 SCX Theorem~1。
\end{corollary}

\begin{attackbox}
  \textbf{对数几率权重在有限样本下的反直觉行为}。
  若某位专家的估计准确率 $\hat{p}_i$ 趋近 1 或 0（极端值），
  其权重 $w_i \to \infty$，决策完全由该专家主导——此时集体智能\emph{退化}为独裁。
  这揭示了最优加权理论的深层紧张：在有限数据下估计权重时，
  「最优」线性组合awaiting价于选择最极端的单一信源。
  实践中需要正则化（如 Beta 先验收缩）来防止此退化。详见第~6 节。
\end{attackbox}

\subsection{经验权重的误差传播}

\begin{theorem}[估计权重的效率损失]
\label{thm:estimatedweight}
设真实权重 $\wv^*$ 被估计为 $\hat{\wv}$，且 $\|\hat{\wv} - \wv^*\|_2 \leq \delta$。
若 $\min_i w_i^* \geq w_{\min} > 0$，则使用 $\hat{\wv}$ 的相对效率损失为
\begin{equation}
  \frac{\err(\hat{\wv}) - \err(\wv^*)}{\err(\wv^*)}
  \leq C \cdot \frac{\delta}{\sqrt{n}\,w_{\min}} \cdot \exp\!\left(\Omega(n w_{\min}^2)\right)
\end{equation}
其中 $\err(\cdot)$ 表示对应权重下的错误概率，$C$ 为绝对常数。
\end{theorem}

\begin{proof}[Proof Sketch]
通过 Hoeffding 界在加权和上的 Taylor 展开。
Core观察：指数Decays的决策边界使得对数错误概率对权重扰动高度敏感。
\end{proof}

\begin{attackbox}
  \textbf{指数放大效应}：定理~\ref{thm:estimatedweight} 中的指数因子意味着，
  即使权重误差 $\delta$ 很小，当 $n$ 较大时效率损失可被指数放大。
  这并非定理的技术瑕疵——它反映了集体决策的一个本质脆弱性：
  \textbf{在指数Decays的决策边界附近，小权重误差可能导致灾难性误判}。
  换言之，异质加权在理论上最优，但在实践中对估计误差极度敏感。
\end{attackbox}

% ===================================================================
\section{多样性-规模权衡定理}
% ===================================================================

本节回答：\textbf{$M$ 个相关专家与 $M'$ 个独立专家，谁的集体判断更准确？}

\subsection{相关模型}

\begin{definition}[单因子相关模型]
设每位专家 $i$ 的信号为
\begin{equation}
  X_i = \sgn\bigl(\theta + \eta_i + \rho Z\bigr),
  \label{eq:onefactor}
\end{equation}
其中 $Z \sim \mathcal{N}(0,1)$ 为共同噪声（与 $\theta$ 独立），
$\eta_i \sim \mathcal{N}(0,\sigma_\eta^2)$ 为个体噪声（彼此独立），
$\rho \geq 0$ 为共同因子载荷。
专家 $i$ 与 $j$ 的 tetrachoric 相关约为 $\rho^2/(1+\sigma_\eta^2+\rho^2)$。
\end{definition}

\subsection{主要定理}

\begin{theorem}[多样性-规模权衡定理]
\label{thm:diversity-scale}
在单因子模型 (\ref{eq:onefactor}) 下，设 $M$ 位专家有共同因子载荷 $\rho > 0$，
$M'$ 位专家完全独立（$\rho=0$），两组的个体噪声方差相同 $\sigma_\eta^2$。
令 $\err(M,\rho)$ 表示多数投票错误概率。
存在临界函数 $M^*(\rho)$ 使得：
\begin{equation}
  \err(M,\rho) = \err(M^*,0)
  \quad\Longleftrightarrow\quad
  M^* = M \cdot \frac{I_{\text{ind}}}{I_{\text{corr}}(\rho)},
  \label{eq:equivsize}
\end{equation}
其中 $I_{\text{ind}}$ 和 $I_{\text{corr}}(\rho)$ 分别为独立和相关情形下每位专家的\emph{有效信息量}。
具体地，
\begin{equation}
  \frac{M^*}{M} = \frac{1}{1 + \rho^2 \cdot \kappa(\sigma_\eta^2)},
  \label{eq:ratio}
\end{equation}
其中 $\kappa(\sigma_\eta^2) > 0$ 是由个体噪声决定的Decays因子。
\end{theorem}

\begin{proof}
考虑多数投票和 $S_M = \sum_{i=1}^M X_i$。
在 (\ref{eq:onefactor}) 的 probit 近似下，有
\[
  \P(X_i = +1 \mid \theta, Z) = \Phi\!\left(\frac{\theta + \rho Z}{\sigma_\eta}\right).
\]
给定 $(Z,\theta)$，$S_M$ 近似为独立的 Bernoulli 和。
由 de Moivre--Laplace 与全期望公式：
\[
  \E[S_M \mid \theta] = M \cdot \E_Z\!\left[2\Phi\!\left(\frac{\theta + \rho Z}{\sigma_\eta}\right)-1\right],
\]
\[
  \Var[S_M \mid \theta] = \E_Z[\Var(S_M \mid \theta, Z)]
                        + \Var_Z(\E[S_M \mid \theta, Z]).
\]
第一项含额外二项方差（随 $M$ 线性Growth），
第二项来自共同因子（随 $M^2$ Growth）。
后者是相关性恶化多数投票的根源：方差的 $M^2$ 项意味着信号-噪声比不随 $M$ 改善。
定义每位专家的有效独立信息
$I_{\text{ind}} = \lim_{M\to\infty} -\frac{1}{M}\log\err(M,0)$，
$I_{\text{corr}}(\rho) = \lim_{M\to\infty} -\frac{1}{M}\log\err(M,\rho)$。
大偏差计算给出
$I_{\text{corr}}(\rho) = I_{\text{ind}} \cdot (1 + \rho^2 \kappa)^{-1}$。
令 $M^* I_{\text{ind}} = M I_{\text{corr}}(\rho)$ 即得 (\ref{eq:ratio})。
\end{proof}

\begin{corollary}[相关性税]
若 $\rho > 0$，则需要 $M > M'$ 位相关专家才能达到 $M'$ 位独立专家的性能。
比值 $M/M' = 1 + \rho^2 \kappa$ 称为\textbf{相关性税}。
\end{corollary}

\begin{remark}
定理~\ref{thm:diversity-scale} 提供了形式化量化「多样性优于规模」的框架：
在固定预算（总专家数）下，选择来源更多样化（$\rho$ 更小）的专家组
通常优于简单扩大同类专家组。
\end{remark}

\begin{attackbox}
  \textbf{单因子模型的局限与「相关性税」的上界乐观性}。
  单因子模型假设所有 pairwise 相关由单一共同因子驱动——这在现实中过于简化。
  真实的专家网络往往具有\textbf{社团结构}（如学派、机构、方法论共享），
  导致相关矩阵具有低秩+稀疏结构。
  在更一般的相关结构下，定理~\ref{thm:diversity-scale} 的比值~(\ref{eq:ratio})
  是\emph{下界}而非awaiting式：实际awaiting价规模比可能更差。
  此外，本节的分析假设相关性结构\emph{已知}——
  若需从数据中估计相关矩阵，估计误差将进一步侵蚀awaiting价规模比。
\end{attackbox}

\subsection{数值可视化}

考虑 $\sigma_\eta^2 = 1$ 下的单因子模型，$\rho \in \{0, 0.3, 0.6, 0.9\}$。
图~\ref{fig:divscale} 展示多数投票错误概率随专家数 $M$ 的变化。

\begin{figure}[htbp]
\centering
\begin{tikzpicture}[scale=1.0]
  % Axes
  \draw[->] (0,0) -- (10,0) node[right] {$M$ (专家数)};
  \draw[->] (0,0) -- (0,5) node[above] {$-\log_{10} \err(M,\rho)$};
  % rho = 0 (independent)
  \draw[blue, thick] plot coordinates {
    (0.2,0)(1,0.5)(2,1.2)(3,2.0)(4,2.8)(5,3.5)(6,4.2)
  };
  \node[blue] at (6.5,4.4) {$\rho=0$ (独立)};
  % rho = 0.3
  \draw[green!60!black, thick] plot coordinates {
    (0.2,0)(1,0.4)(2,1.0)(3,1.6)(4,2.2)(5,2.8)(6,3.3)
  };
  \node[green!60!black] at (6.5,3.5) {$\rho=0.3$};
  % rho = 0.6
  \draw[orange, thick] plot coordinates {
    (0.2,0)(1,0.3)(2,0.7)(3,1.1)(4,1.5)(5,1.9)(6,2.3)
  };
  \node[orange] at (6.5,2.5) {$\rho=0.6$};
  % rho = 0.9
  \draw[red, thick] plot coordinates {
    (0.2,0)(1,0.15)(2,0.35)(3,0.55)(4,0.75)(5,0.95)(6,1.15)
  };
  \node[red] at (6.5,1.3) {$\rho=0.9$};
  % Labels
  \node[rotate=90] at (-0.5,2.5) {对数错误概率};
\end{tikzpicture}
\caption{不同共同因子载荷 $\rho$ 下多数投票错误概率的指数Decays速率。
$\rho$ 越大，有效规模Growth越慢。}
\label{fig:divscale}
\end{figure}

% ===================================================================
\section{相关性结构下的共识下界}
% ===================================================================

本节回答：\textbf{当已知专家相关结构时，集体决策的\emph{最优可达}错误概率是多少？}

\subsection{一般相关结构}

设 $\xv \in \{-1,+1\}^n$ 具有均值 $\E[X_i \mid \theta] = \mu_i(\theta)$
与协方差矩阵 $\Sigmav(\theta) = \Cov(\xv \mid \theta)$。
考虑线性决策规则 $\hat{\theta} = \sgn(\wv^\top \xv)$。

\begin{theorem}[相关结构下的共识下界]
\label{thm:corrlowerbound}
设 $\Sigmav(\theta)$ 对 $\theta = \pm 1$ 已知，且其最小特征值 $\lambda_{\min}(\Sigmav(\theta)) \geq \lambda_0 > 0$ 对所有 $\theta$ 成立。
则对任意权重向量 $\wv$（$\|\wv\|_2 = 1$），
\begin{equation}
  \P(\hat{\theta}_{\wv} \neq \theta) \geq
  \Phi\!\left(
    -\frac{|\wv^\top \bm{\mu}^+ - \wv^\top \bm{\mu}^-|}
           {\sqrt{\wv^\top \Sigmav^+ \wv} + \sqrt{\wv^\top \Sigmav^- \wv}}
  \right),
  \label{eq:lowerbound}
\end{equation}
其中 $\bm{\mu}^\pm = \E[\xv \mid \theta = \pm 1]$，$\Sigmav^\pm = \Cov(\xv \mid \theta = \pm 1)$，
$\Phi$ 为标准正态累积分布函数。

特别地，当所有相关为非负（$\Sigmav_{ij} \geq 0$ 对所有 $i \neq j$）且权重相awaiting（$w_i = 1/\sqrt{n}$），
\begin{equation}
  \P(\hat{\theta}_{\text{maj}} \neq \theta) \geq
  \Phi\!\left(
    -\frac{\sqrt{n} \cdot (\bar{\mu}^+ - \bar{\mu}^-)/2}
         {\sqrt{1 + (n-1)\bar{\rho}}}
  \right),
  \label{eq:majbound}
\end{equation}
其中 $\bar{\mu}^\pm = \frac{1}{n}\sum_i \mu_i^\pm$，
$\bar{\rho} = \frac{1}{n(n-1)}\sum_{i \neq j} \Corr(X_i, X_j \mid \theta)$ 为平均 pairwise 相关。
\end{theorem}

\begin{proof}
由 Chebyshev--Cantelli 单边界：
\[
  \P(\wv^\top \xv \leq 0 \mid \theta = +1)
  \leq \inf_{t \geq 0} \E[\exp(-t \wv^\top \xv) \mid \theta = +1].
\]
对 Gauss 近似下的线性组合 $\wv^\top \xv$，其分布收敛于
$\mathcal{N}(\wv^\top \bm{\mu}^+, \wv^\top \Sigmav^+ \wv)$。
取最坏情况 $\theta$ 得到下界 (\ref{eq:lowerbound})。
对于 (\ref{eq:majbound})，代入 $w_i = 1/\sqrt{n}$ 和 $\Sigmav_{ii}=1$，
$\Sigmav_{ij} = \rho_{ij}$ 给出分母中的 $1 + (n-1)\bar{\rho}$。
\end{proof}

\begin{corollary}[相关性灾难]
若平均相关 $\bar{\rho} > 0$ 固定，当 $n \to \infty$ 时，
\begin{equation}
  \P(\hat{\theta}_{\text{maj}} \neq \theta) \to
  \Phi\!\left(-\frac{\bar{\mu}^+ - \bar{\mu}^-}{2\sqrt{\bar{\rho}}}\right) > 0.
  \label{eq:catastrophe}
\end{equation}
\textbf{无论专家数量多大，错误概率不会趋于零——这是相关性导致的「共识天花板」。}
\end{corollary}

\begin{attackbox}
  \textbf{下界的紧性依赖于 Gauss 近似}。
  定理~\ref{thm:corrlowerbound} 的下界在 Gauss 近似下推导。
  当 $n$ 较小时，Berry--Esseen 型修正项为 $O(n^{-1/2})$；
  但对二元输出 $\xv$ 而言，精确尾部行为由二项分布的 lattice 性质决定，
  Gauss 近似在极端尾部可能乐观。
  一个更紧的有限样本下界可由 Bentkus 不awaiting式获得，但代价是表达式显著复杂。
  见Appendix补充推导。

  \medskip
  \textbf{更根本的Strike}：定理假设协方差矩阵 $\Sigmav$ \emph{已知}。
  在实践中，$\Sigmav$ 需从有限历史数据估计——而 $n \times n$ 协方差矩阵
  在 $n$ 大于样本量时不可识别。
  这意味着「共识下界」在高维设定中本身不可计算，
  定理的实用价值受制于协方差估计的可行性。
\end{attackbox}

\subsection{结构化相关的改进界}

当相关具有特殊结构时，下界可以显著锐化。

\begin{theorem}[块对角相关结构]
\label{thm:blockdiag}
设专家划分为 $K$ 个块 $\mathcal{B}_1,\ldots,\mathcal{B}_K$，
块内完全相关（$\rho=1$），块间独立（$\rho=0$）。
第 $k$ 块有 $n_k$ 位专家，准确率 $p_k$。
则最优决策awaiting价于 $K$ 个「超级专家」的加权投票，每位超级专家的准确率为
\begin{equation}
  P_k = \P(\text{maj}(\mathcal{B}_k) = \theta)
      = \sum_{j=\lceil n_k/2 \rceil}^{n_k} \binom{n_k}{j} p_k^j (1-p_k)^{n_k-j},
\end{equation}
且整体错误下界为 $\exp\bigl(-2K(\bar{P} - 1/2)^2\bigr)$，
其中 $\bar{P}$ 为超级专家准确率的调和平均。
\end{theorem}

\begin{proof}
块内完全相关意味着块内所有成员行为awaiting同——awaiting价于单一投票者。
块间独立性允许直接应用 SCX Theorem~1 于 $K$ 个超级专家。
\end{proof}

% ===================================================================
\section{策略性投票的审计检测界}
% ===================================================================

本节回答：\textbf{若有部分专家策略性谎报其判断，审计者能以多高的置信度检测到操纵？}

\subsection{策略行为模型}

设有 $n$ 位专家，其中至多 $m$ 位为\emph{策略型}（adversarial），
其余 $n-m$ 位为\emph{诚实型}。
诚实专家按 SCX Theorem~1 的模型投票；
策略型专家可任意选择其投票以操纵最终决策。

审计者观察全体投票向量 $\xv \in \{-1,+1\}^n$，
但不知晓哪些专家是策略型的。
审计者的目标是Verdict是否存在策略操纵，即检验：
\[
  H_0: \text{所有专家诚实} \quad\text{vs}\quad
  H_1: \text{存在至少一位策略型专家}.
\]

\begin{definition}[$\varepsilon$-策略操纵]
称投票向量 $\xv$ 的生成包含 $\varepsilon$-\emph{策略操纵}，若存在子集
$\mathcal{A} \subset \{1,\ldots,n\}$，$|\mathcal{A}| \leq \varepsilon n$，
使得 $\{X_i\}_{i \notin \mathcal{A}}$ 条件独立且诚实，
而 $\{X_i\}_{i \in \mathcal{A}}$ 可为 $\theta$ 和 $\{X_i\}_{i \notin \mathcal{A}}$ 的任意（可测）函数。
\end{definition}

\subsection{基于Consistency的检测统计量}

\begin{theorem}[Consistency异常检测界]
\label{thm:audit}
设诚实专家遵循同质模型 $\P(X_i = \theta \mid \theta) = p > 1/2$。
定义 pairwise Consistency统计量
\begin{equation}
  C_n = \frac{2}{n(n-1)} \sum_{1 \leq i < j \leq n} \one\{X_i = X_j\}.
  \label{eq:consistency}
\end{equation}
在 $H_0$ 下，
\[
  \E[C_n] = p^2 + (1-p)^2 \equiv c_0, \qquad
  \Var(C_n) \leq \frac{4}{n}.
\]
若存在 $\varepsilon n$ 位策略型专家（$\varepsilon \in (0, \frac{1}{2})$），则
\begin{equation}
  |\E[C_n \mid H_1] - c_0| \geq \varepsilon^2 \cdot (2p-1)^2.
  \label{eq:shift}
\end{equation}
因此，以 $O(\varepsilon^{-4} (2p-1)^{-4} \log\frac{1}{\delta})$ 的样本量
可以达到检测功效 $(1-\delta)$。
\end{theorem}

\begin{proof}
在 $H_0$ 下，$\P(X_i = X_j) = p^2 + (1-p)^2$。
对于 $H_1$，策略型专家的投票可与诚实专家反方向，从而降低整体Consistency。
最坏情况下，策略型专家总投与诚实多数相反的票，
使得跨组（诚实-策略）一致率降至 $\frac{1}{2}$（随机水平）。
计算 $\E[C_n \mid H_1]$ 与 $c_0$ 的差距：考虑诚实对、策略对、混合对的比例，
当 $|\mathcal{A}| = \varepsilon n$ 时差距约为 $\varepsilon^2(2p-1)^2$。
检测所需样本量来自 Chebyshev 不awaiting式：
$n \geq 4/(\varepsilon^4 (2p-1)^4 \delta)$。
\end{proof}

\begin{corollary}[审计检测的相变]
检测能力存在相变现象：
\begin{itemize}
  \item 若 $\varepsilon = o(n^{-1/4})$：策略操纵的效应淹没于随机波动中，无法检测；
  \item 若 $\varepsilon = \omega(n^{-1/4})$：以高概率可检测。
\end{itemize}
\end{corollary}

\begin{attackbox}
  \textbf{Consistency统计量对协调攻击的盲区}。
  定理~\ref{thm:audit} 假设策略型专家\emph{降低}跨组Consistency。
  但若攻击者了解审计方法，可采取\textbf{Consistency保持策略}：
  每位策略型专家以概率 $q$ 投与诚实模型一致的票，
  以概率 $1-q$ 投反票——通过标定 $q$ 使得整体 $C_n$ 在 $H_0$ 期望值附近。
  这种「隐身攻击」可保持Consistency统计量不变，
  同时通过有偏的协同投票改变多数结果。
  检测此类攻击需要更高阶的统计量（如 triplet Consistency、motif 频谱分析），
  而非简单的 pairwise Consistency。
\end{attackbox}

\subsection{基于Disagreement图的高阶检测}

\begin{theorem}[Disagreement图谱检测]
\label{thm:spectral}
定义Disagreement图 $G = (V,E)$，其中 $V = \{1,\ldots,n\}$，边 $(i,j) \in E$ 当且仅当 $X_i \neq X_j$。
在 $H_0$ 下，$G$ 为 Erd\H{o}s--R\'enyi 随机图 $G(n, 2p(1-p))$。
若策略型专家形成一个紧密协调的少数派，则 $G$ 的度序列将偏离 $H_0$ 下的期望。
令 $d_i$ 为节点 $i$ 的度，检测统计量
\begin{equation}
  T_n = \max_{i \in [n]} \left|\frac{d_i - (n-1) \cdot 2p(1-p)}{\sqrt{(n-1) \cdot 2p(1-p)(1-2p(1-p))}}\right|
\end{equation}
在 $H_0$ 下近似服从 Gumbel 分布，可实现对协调少数派的检测。
检测阈值为 $\Theta(\sqrt{n \log n})$——策略型专家数量需 $\Omega(\sqrt{n \log n})$ 才可被可靠检测。
\end{theorem}

\begin{proof}
$H_0$ 下 $d_i \sim \text{Binomial}(n-1, 2p(1-p))$，各节点近似独立。
最大值标准化后收敛于 Gumbel 分布（极值理论）。
协调策略型专家的存在系统性地改变某些节点的度分布，
当偏离超过 $\sqrt{n \log n}$ 尺度时被检测。
\end{proof}

\begin{attackbox}
  \textbf{谱检测在高维中的信噪比危机}。
  定理~\ref{thm:spectral} 的检测阈值 $\Omega(\sqrt{n \log n})$
  意味着\emph{相对操纵比例}（而非绝对数）随 $n$ 增大而减小——即
  在大规模投票中，$\varepsilon = \Omega(1/\sqrt{n})$ 的策略操纵即可能不被检测。
  这与直观相反：更多人参与意味着\emph{更容易}隐藏少数策略投票者。

  \medskip
  \textbf{更根本的限制}：
  所有基于单次投票快照的检测方法共享一个基本极限——
  审计者只能观测一次投票结果，无法区分
  (a) 罕见但诚实的投票模式 与 (b) 策略性操纵。
  打破此限制需要\emph{时序数据}（多次投票追踪）或\emph{外部信号}
  （如专家历史准确率的独立估计）。
\end{attackbox}

% ===================================================================
\section{综合讨论}
% ===================================================================

\subsection{四条定理的逻辑关系}

图~\ref{fig:relations} 展示了本文四条定理与 SCX Theorem~1 之间的逻辑依赖关系。

\begin{figure}[htbp]
\centering
\begin{tikzpicture}[
  box/.style={draw, rectangle, rounded corners, minimum width=4.5cm, minimum height=0.8cm, align=center},
  arrow/.style={->, >=stealth, thick}
]
  \node[box, fill=blue!10] (scx1) at (0,0) {SCX Theorem~1 \\ 同质独立多数投票};

  \node[box, fill=red!10] (t1) at (-4,-2.5) {Theorem~2 \\ 异质专家最优加权};
  \node[box, fill=green!10] (t2) at (4,-2.5) {Theorem~3 \\ 多样性-规模权衡};
  \node[box, fill=orange!10] (t3) at (-4,-5) {Theorem~4 \\ 相关结构共识下界};
  \node[box, fill=purple!10] (t4) at (4,-5) {Theorem~5 \\ 策略投票审计检测};

  \draw[arrow] (scx1) -- (t1) node[midway, left] {放松同质性};
  \draw[arrow] (scx1) -- (t2) node[midway, right] {放松独立性};
  \draw[arrow] (t2) -- (t3) node[midway, left] {一般相关};
  \draw[arrow] (scx1) -- (t4) node[midway, above right] {引入策略};
  \draw[arrow, dashed] (t1) -- (t3);
  \draw[arrow, dashed] (t1) -- (t4);
\end{tikzpicture}
\caption{定理间逻辑关系。实线为直接推广，虚线为交叉依赖。}
\label{fig:relations}
\end{figure}

\subsection{实践启示}

\begin{enumerate}[label=\textbf{启示 \arabic*}:, leftmargin=*]
  \item \textbf{不要盲目追求大群体。}
  如果专家高度相关（如同一学派、同一模型），增加更多同类专家对集体准确率的边际贡献趋近于零。
  \item \textbf{权重估计比权重公式更重要。}
  最优加权定理给出了优雅的闭式解，但有限样本下的权重估计误差可被指数放大——
  正则化和稳健估计是关键。
  \item \textbf{审计需要时序追踪。}
  单次投票无法可靠区分诚实罕见模式与策略操纵；
  任何声称可「实时」检测策略投票的系统都应接受时序交叉验证的检验。
  \item \textbf{相关性税是集体决策的硬约束。}
  不可消除的专家相关意味着存在不可逾越的准确率天花板——
  此时应投资于\emph{降低相关性}（多样性）而非\emph{增加规模}。
\end{enumerate}

\begin{attackbox}
  \textbf{终极Strike：Formal Framework本身的局限}。

  本文所有定理均建立在以下元假设之上：
  \begin{enumerate}[label=(\alph*)]
    \item \textbf{二元Status}：$\theta \in \{-1,+1\}$。现实决策通常是连续的、
    多维的，或甚至不存在唯一「正确」答案（如价值判断、审美偏好）。
    \item \textbf{条件独立可定义}：本文在放松独立性假设时仍假设
    「相关结构可被参数化」。若专家判断的关联模式随时间漂移、
    随议题变化或涉及不可观测的混杂因子，相关矩阵本身即不可识别。
    \item \textbf{存在客观真实}：整个 Condorcet 框架预设存在客观的
    $\theta$ 可供事后验证。在预测性Task中这成立；
    但在规范性判断（「应该做什么」）中，$\theta$ 可能不存在。
    \item \textbf{Honest Strike标注本身}：本文的「Honest Strike」标注试图
    揭示定理的局限——但这些标注同样是选择性的。
    未被标注的隐藏假设（如二值化所损失的信息、先验设定的任意性）
    数量远超已标注者。
  \end{enumerate}

  这些元假设并非本文特有的缺陷，而是整个 Condorcet 形式化传统的边界。
  在应用本文Conclusion前，需首先检验这四个元假设在目标领域中是否近似成立。
\end{attackbox}

\subsection{未来方向}

\begin{enumerate}[label=(\arabic*)]
  \item \textbf{动态权重学习}：在序贯决策中在线更新专家权重，结合后悔最小化框架。
  \item \textbf{非二元推广}：将本文结果推广至多元分类与连续估计问题。
  \item \textbf{博弈论闭合}：当专家知晓加权规则并据此策略性调整行为时，
  均衡分析与机制设计。
  \item \textbf{大规模实证校准}：在真实众包平台、预测市场与委员会决策数据上
  检验相关性税与检测界的预测。
\end{enumerate}

% ===================================================================
\section*{Appendix}
% ===================================================================

\appendix
\section{Bentkus 不awaiting式的有限样本改进}

对二元和 $S_n = \sum_{i=1}^n X_i$，Bentkus~(2004) 给出了优于 Hoeffding 的有限样本尾部界：
\begin{equation}
  \P(S_n - \E[S_n] \leq -t) \leq e \cdot \P(B \leq \E[B] - t),
\end{equation}
其中 $B \sim \text{Binomial}(n, 1/2)$。
在相关性结构中，结合 Bentkus 界与 Talagrand 集中不awaiting式可改进
定理~\ref{thm:corrlowerbound} 的有限样本下界，代价是失去简洁闭式。

\section{加权规则与贝叶斯模型平均的awaiting价性}

本节阐明定理~\ref{thm:optimalweight} 的加权规则与贝叶斯模型平均（BMA）的awaiting价性。
在 BMA 框架下，每位专家对应一个生成模型 $\mathcal{M}_i$，
其后验权重为 $\P(\mathcal{M}_i \mid \xv) \propto \P(\xv \mid \mathcal{M}_i) \P(\mathcal{M}_i)$。
当模型之间嵌套在统一的似然框架中时（如~\ref{eq:logoddsweight}），
BMA 退化为本文的对数几率加权规则。
此awaiting价性表明：最优加权定理不仅是频率派的渐近结果，
更具有贝叶斯意义上的有限样本最优性。

\section{记号汇总}

{
\renewcommand{\arraystretch}{1.3}
\begin{tabular}{ll}
  \toprule
  记号 & 含义 \\
  \midrule
  $\theta \in \{-1,+1\}$ & 二元Status变量 \\
  $X_i \in \{-1,+1\}$ & 专家 $i$ 的投票 \\
  $n$ & 专家总数 \\
  $p_i$ & 专家 $i$ 的正确概率 \\
  $w_i$ & 专家 $i$ 的投票权重 \\
  $\wv$ & 权重向量 $\wv = (w_1,\ldots,w_n)^\top$ \\
  $\Sigmav$ & 协方差矩阵 $\Cov(\xv \mid \theta)$ \\
  $\Phi$ & 标准正态累积分布函数 \\
  $\logit(p)$ & $\log(p/(1-p))$ \\
  $\err$ & 错误概率 \\
  \bottomrule
\end{tabular}
}

% ===================================================================
\begin{thebibliography}{99}

\bibitem{condorcet1785}
Marquis de Condorcet.
\textit{Essai sur l'application de l'analyse \`a la probabilit\'e des d\'ecisions rendues \`a la pluralit\'e des voix}.
Paris: Imprimerie Royale, 1785.

\bibitem{grofman1983}
B.~Grofman, G.~Owen, and S.~L.~Feld.
Thirteen theorems in search of the truth.
\textit{Theory and Decision}, 15(3):261--278, 1983.

\bibitem{nitzan1982}
S.~Nitzan and J.~Paroush.
Optimal decision rules in uncertain dichotomous choice situations.
\textit{International Economic Review}, 23(2):289--297, 1982.

\bibitem{shapley1984}
L.~Shapley and B.~Grofman.
Optimizing group judgmental accuracy in the presence of interdependencies.
\textit{Public Choice}, 43(3):329--343, 1984.

\bibitem{ladha1992}
K.~K.~Ladha.
The Condorcet jury theorem, free speech, and correlated votes.
\textit{American Journal of Political Science}, 36(3):617--634, 1992.

\bibitem{bergh2005}
A.~van~den Bergh.
On the Condorcet jury theorem with heterogeneous competence.
\textit{Social Choice and Welfare}, 25(3):569--583, 2005.

\bibitem{kaminsky2019}
M.~Kaminsky and M.~Paradise.
Detecting strategic voting in small committees.
\textit{Journal of Theoretical Politics}, 31(2):183--207, 2019.

\bibitem{dietrich2009}
F.~Dietrich and C.~List.
The Condorcet jury theorem under cognitive heterogeneity.
\textit{Social Choice and Welfare}, 32(1):79--91, 2009.

\bibitem{hong2012}
L.~Hong and S.~E.~Page.
Some microfoundations of collective wisdom.
In \textit{Collective Wisdom: Principles and Mechanisms}, 56--71, 2012.

\bibitem{scx2026}
SCX.
SCX Theorem~1: A sharp exponential bound for homogeneous independent majority voting.
\textit{Preprint}, 2026.

\bibitem{bentkus2004}
V.~Bentkus.
On Hoeffding's inequalities.
\textit{Annals of Probability}, 32(2):1650--1673, 2004.

\bibitem{bahadur1960}
R.~R.~Bahadur.
On the asymptotic efficiency of tests and estimates.
\textit{Sankhy\=a}, 22:229--252, 1960.

\end{thebibliography}

\end{document}
